\chapter*{Abstract}
\addcontentsline{toc}{chapter}{Abstract}

Network intrusion detection systems (NIDS) increasingly rely on supervised
machine learning trained on flow-level features. Public benchmarks such as
CICIDS2017 and UNSW-NB15 routinely report near-perfect macro F1 scores, yet
it remains unclear whether these numbers reflect operational robustness or
in-distribution memorisation. This thesis closes that gap empirically. The
originality is not in the choice of models --- random forest, XGBoost, and
LightGBM are standard --- but in the combination of honest temporal
evaluation, score-query black-box adversarial robustness analysis, and
SOC-grade MITRE ATT\&CK triage in a single reproducible pipeline.

Four classical models --- logistic regression, random forest, XGBoost, and
LightGBM --- are trained on CICIDS2017 under three evaluation protocols:
random stratified, day-based temporal, and leave-one-day-out (LODO).
UNSW-NB15 is used as a cross-dataset validation set under its official
train/test split. Both multi-class and binary tasks are evaluated. We then
probe robustness with a score-query black-box greedy evasion attack on the
binary day-split random forest, quantify cross-model transferability, and
test naïve adversarial training as a defence. Finally, predicted attack
types are mapped to MITRE ATT\&CK techniques to produce SOC-actionable
alerts.

The headline findings are fourfold. First, multi-class macro F1 collapses
from 0.86 (stratified XGBoost) to 0.44 on the day split, with all four
models within 0.04 of one another, because Friday's attack classes never
appear in training. Second, binary detection survives the temporal shift:
random forest reaches F1 = 0.86 on the day split and ROC-AUC = $0.93 \pm
0.05$ across LODO folds. Third, the random forest is brittle under
low-effort evasion (21.4\% of attack flows flipped at $\varepsilon =
0.25\sigma$, with a looser headline evasion rate of 24.4\% when the attack
is not capped, on a 19-feature subset), and adversarial examples transfer
to XGBoost at 86.4\%. Fourth, naïve adversarial training does not help ---
clean F1 drops from 0.86 to 0.73 while evasion stays at 26.8\%.

We conclude that headline macro-F1 on stratified splits substantially
overstates deployment readiness under these benchmarks, and recommend
day-based or LODO splits as the default evaluation protocol for flow-based
intrusion detection research.

\bigskip
\noindent\textbf{Keywords:} network intrusion detection; flow-based features; CICIDS2017; UNSW-NB15; random forest; XGBoost; LightGBM; temporal generalisation; adversarial robustness; MITRE ATT\&CK.
